\documentclass[11pt,a4paper]{article}
\usepackage[utf8]{inputenc}
\usepackage[T1]{fontenc}
\usepackage[spanish]{babel}
\usepackage{hyperref}
\usepackage{geometry}
\geometry{margin=2.5cm}
\usepackage{parskip}

\title{091 el absoluto no se combate al final- se impide al enlazar}
\author{Alejandro Valenzuela Rivera\\\small Universidad Autónoma de Nuevo León}
\date{\today}
\begin{document}
\maketitle

\begin{quote}
\textbf{Sesgo declarado.} Este trabajo se sostiene desde el sesgo declarado de su autor, y no finge neutralidad. Lo que afirma se puede tumbar, y ese es el punto: el falsificador está escrito, no escondido.
\end{quote}

\section*{Resumen}
Nada existe si no esta enlazado a una pieza de donde viene, y lo que no se busco no es un vertice: es una ausencia. El absoluto no se corrige despues de escrito; se impide en el momento en que se escribe, porque ahi es cuando se puede decir de donde viene.

\section{Por qué se sostiene}
Y hay seis puertas conocidas por las que se cuela, una por cada forma en que se disfraza: el nombre de marca, la cantidad sin fuente, la fuente sin enlace, la utilidad no medida, el silencio tomado por vertice, y la regla que no declara enlace.

\section{De dónde viene}
La autoridad y el repo, 2026-10-25.

\section{El caso}
Ninguna: es la regla del sistema.

\section{Método}
Una pieza que se puede contestar antes de que se vuelva regla.

\section{Límites, y qué lo tumba}
Una regla que se impone y solo se puede discutir cuando ya hizo dano.

\section{Lo que queda abierto}
\textquestiondown{}Si el absoluto se impide al enlazar, que pasa con lo que quedo enlazado mal?

\section*{Cómo se leyó esta tesis}
Esto se escribió en el repositorio man-made-empathy, del que este candado es una pieza: candados/04 el lenguaje/091 el absoluto no se combate al final- se impide al enlazar.md. La máquina que lo compuso lo hizo desde este mismo archivo, y no agregó nada que no estuviera escrito. Si algo de esto es falso, se dice en el mismo sitio, y se corrige con fecha.

\end{document}
